% Einheit 5 von 5 – Stationär und langfristig (bank/matrizen-und-uebergangsprozesse/e5.jsonl, zusammenbau v0.1)
%% TODO Zeitmarke Einheit 5: Marken-Zeile nicht lesbar
%% TODO Zweigzeile Teil 1: was hier gelernt wird (Satz in Schülersprache aus dem Einheitstitel) – Einheit 5
\einheitenkopf[e5]{Einheit 5 von 5 · Stationär und langfristig}
\zweigzeile{keine Prüfungsaufgabe}
\setcounter{aufgabe}{16}

%% TODO Titel: Ich-kann-Satz fehlt in der Bank (Platzhalter: Kettenname) – Nr. 17
%% TODO Anweisung: ein Satz über der ersten Teilaufgabe fehlt in der Bank – Nr. 17
\begin{aufgabe}{Stationär und langfristig}
%% TODO \rechenplatz{n}: Schrittzahl des Grundfalls fehlt in der Bank (nur bei fester Schreibform, 2.3 b)
\begin{teile}
\teil Es gilt $v_{n+1} = M \cdot v_n$; v ist die heutige Verteilung. Was beschreibt $M^3 \cdot v$? \\ \kreuz{drei Schritte vorwärts} \\ \kreuz{drei Schritte rückwärts} \\ \kreuz{einen Schritt vorwärts}
\end{teile}
\begin{gleichungsraster}[2]
\gl{\text{Es gilt $v_{n+1} = M \cdot v_n$; $M = ((0{,}8 | 0{,}4), (0{,}2 | 0{,}6))$, Gesamtzahl 900. Berechne die stationäre Verteilung.}} &
\end{gleichungsraster}
\begin{teile}
\teil Es gilt $v_{n+1} = M \cdot v_n$; $M = ((0{,}6 | 0), (0{,}4 | 1))$ (Reihenfolge A, B), Gesamtzahl 500. Gib die stationäre Verteilung an. ( \leerfeld | \leerfeld )
\end{teile}
\begin{gleichungsraster}[2]
\gl{\text{Es gilt $v_{n+1} = M \cdot v_n$; $M = ((0{,}8 | p), (0{,}2 | 1 - p))$; die Verteilung $(300 | 200)$ ist stationär. Bestimme p.}} &
\end{gleichungsraster}
\begin{teile}
\teil Es gilt $v_{n+1} = M \cdot v_n$; $M = ((0{,}9 | 0{,}2), (0{,}1 | 0{,}8))$ (Reihenfolge A, B), $v = (800 | 400)$. Wie viele wechseln von A nach B, wie viele von B nach A? Was folgt daraus?
\teil Es gilt $v_{n+1} = M \cdot v_n$; $M = ((0 | 0 | 4), (0{,}5 | 0 | 0), (0 | 0{,}5 | 0))$. Zeige $M^3 = E$ und deute das Ergebnis.
\teil Es gilt $v_{n+1} = M \cdot v_n$; $M = ((0 | 3), (0{,}5 | 0{,}5))$ und $v = (6 | 3)$. Zeige, dass $M \cdot v$ ein Vielfaches von v ist, und gib den Faktor an.
\teil Eine Lichterkette wechselt jede Sekunde zwischen Rot, Gelb und Blau (in dieser Reihenfolge) nach $K_b = ((b | 2b | 0), (2b | b | 0), (1 - 3b | 1 - 3b | 1))$ mit $0 \le b \le \frac{1}{3}$. Für $b < \frac{1}{3}$ nähern sich die Potenzen $K_b^k$ der Matrix mit den Zeilen $(0 | 0 | 0)$, $(0 | 0 | 0)$, $(1 | 1 | 1)$. Beschreibe die Farbwechsel in Abhängigkeit von b \hfill (Abitur 2022 LK).
\end{teile}
\end{aufgabe}

%% TODO Titel: Ich-kann-Satz fehlt in der Bank (Platzhalter: Kettenname) – Nr. 18
%% TODO Anweisung: ein Satz über der ersten Teilaufgabe fehlt in der Bank – Nr. 18
\begin{aufgabe}{Übergangsprozess: Eignung eines Populationsmodells zur langfristigen Beschreibung beurteilen}
\begin{teile}
\teil Eine Vogelpopulation auf einer kleinen Insel wächst im Modell jedes Jahr um den Faktor $1{,}2$. Beurteile, ob das Modell die Entwicklung über 50 Jahre beschreiben kann.
\end{teile}
\end{aufgabe}

%% TODO Titel: Ich-kann-Satz fehlt in der Bank (Platzhalter: Kettenname) – Nr. 19
%% TODO Anweisung: ein Satz über der ersten Teilaufgabe fehlt in der Bank – Nr. 19
\begin{aufgabe}{Stationär und langfristig – Fehler finden, Begründen, Anwendung}
\begin{teile}
\teil Es gilt $v_{n+1} = M \cdot v_n$; $M = ((0{,}9 | 0{,}4), (0{,}1 | 0{,}6))$, der Verein hat 1000 Mitglieder in zwei Gruppen. Anna sagt: „Stationär ist $(4 | 1)$.“ Finde den Fehler und gib die stationäre Verteilung an.
\teil Begründe, warum man beim Fixvektor einer Übergangsmatrix mit zwei Zuständen eine der beiden Gleichungen durch die Gesamtzahl ersetzt.
\teil In einem Ort gibt es zwei Bäckereien. Von A wechseln jede Woche 10\,\% der Kunden zu B, von B 15\,\% zu A; zusammen sind es 2000 Kunden. Bäckerei A braucht ab 1100 Kunden eine zweite Verkäuferin. Braucht sie auf Dauer eine?
\end{teile}
\end{aufgabe}
